%% kapa: exam=1 round=21 year=2010 period=2 form=A subject=accounting subject_ko=회계학 seq=40 number=40
%% source: Q-Net 감정평가사 기출문제 (cst00309) artlSeq=5024181 file=extracted/1008/5024181/2010년 제21회 감정평가사 1차시험 기출문제/2교시.hwp
%% conversion: kordoc 4.24.51 -> tools/gen1.py (auto); review=proofread-sonnet; flags=none
\begin{question}{40}
(주)한국은 아래 표와 같이 제조부문과 보조부문을 각각 두 개씩 운영하고 있다. 동력부와 관리부에 집계된 제조간접원가는 각각 ￦200,000과 ￦370,000이다. 각 부문간의 용역수수비율은 아래 표와 같다. 보조부문의 제조간접원가를 상호배부법으로 배부할 때, 제과부에 배부되는 제조간접원가는 얼마인가?

{\small\centering
\begin{tabular}{|c|c|c|c|c|}
\hline
\multirow{2}{*}{\makecell{사용처\\제공처}} & \multicolumn{2}{c|}{제조부문} & \multicolumn{2}{c|}{보조부문} \\
\cline{2-5}
 & 제과부 & 제빵부 & 동력부 & 관리부 \\
\hline
동력부 & 0.2 & 0.4 & - & 0.4 \\
\hline
관리부 & 0.4 & 0.35 & 0.25 & - \\
\hline
\end{tabular}\par}

\begin{choices}
  \item ￦265,000
  \item ￦285,000
  \item ￦305,000
  \item ￦325,000
  \item ￦345,000
\end{choices}
\end{question}
\answer{1}
